Wir vergleichen ähnliche Körper (gleiche Form, gleiches Material) mit der typischen Länge $L$. Dann gilt
\begin{align}
 m_\text{s} &\propto V \propto L^3 & m_\text{t} &\propto A \propto L^2
 \quad\Rightarrow\quad \frac{m_\text{s}}{m_\text{t}} \propto L
\end{align}
Das Verhältnis von schwerer zu träger Masse hängt also von der Grösse des Körpers ab. Äquivalent sind die beiden Massen nur, wenn sie für alle Körper proportional zueinander sind. Das ist hier nicht der Fall: $\result{\text{Die Massen sind nicht äquivalent.}}$

Beobachtbar wäre das zum Beispiel beim freien Fall: Die Gewichtskraft ist proportional zur schweren Masse, die Beschleunigung ist Kraft durch träge Masse,
\begin{align}
 a &= \frac{F_\text{G}}{m_\text{t}} \propto \frac{m_\text{s}}{m_\text{t}} \propto L
\end{align}
Grosse Körper würden schneller fallen als kleine.
